\gddchapter{Analysis}

\section{Thematic Analysis Framework}
Six phase model to move from theoretical raw audio to theoretical insights:
This is just a scaffold that will be filled out with actual data from this session:


% \begin{enumerate}
%     \item Familiarisation
    
%     Listen to the recorings while reading chapter 5 to notice intiial patterns

%     \item Initial Coding
    
%     Generate labels for specific UX moments such as "Command frustration", "modal confusion" etc.

%     \item Searching for Themes
    
%     Group codes into broader categories like "Neutral language friction" or "Immersion breaking"

%     \item Reviewing Themes
    
%     Check if those found themes accurately represent the difference between two input modalities, 
%     and if they fit the main research question.

%     \item Defining / Refining Themes
    
%     Refine the essence of what the data tells you about your research themes

%     \item Writing Up
    
%     Integrate transcript extracts in the finial presentation that show how the player 
%     interacted with the experience, including the themes defined above as anchoring points.

% \end{enumerate}


\subsection{Familiarisation}

Some initial patterns that were noted after listening to the recording:

\begin{enumerate}
    \item She registered the "Let's" just like \#2
    \item The participant seems to be goal-oriented
    \item She treated the descriptions very seriously, treating them as part of the puzzle not just navigation guidance
    \item Klara asked twice if she can interact with the objects not listed in the description showcasing the curiosity and boundary-probing.
    \item She attributed feeling overwhelmed to the novelty of the voice-driven experience rather than the modality itself.
    \item Her answers on immerssion are inconsistent pointing to more nuance in her experience.
    \item She framed the voice based navigation as the ChatGPT era and keyboard navigation as 80s Sierra text based games. 
    \item The participant made remarks about seeking secrets hidden in the game and her desire to look for them.
    \item She appreciated the writing style of the descriptions 
    \item Klara was wearing a medical mask during the testing and expressed her concern over the accuracy of voice recognition.
 
    
\end{enumerate}

\subsection{Initial coding}

\begin{enumerate}
    \item Cooperative register
    
        \begin{quote}
            "I feel like we are playing together, me and the computer... it's like we are together in this." 
        \end{quote}
        The experience marked as a cooperatie endavour not a solo playthough. 


    

    \item Illusion of choice
    
        \begin{quote}
            "I think it gives you the illusion of choice, because you don't know what you can do and what you cannot."
        \end{quote}
        Open command environment named as ambiguity rather than a freedom or overwhelming space

        \item First step anxiety
    
        \begin{quote}
            "How to take the first step? How to start? That was my main concern." 
        \end{quote}
        The player seems to be overwhelemed at the beginning despite the brief intro.


        \item AI era vs retro framing
    
        \begin{quote}
            "First version makes me think more of ChatGPT... the other version makes me think of those old 80s games from Sierra."
        \end{quote}
        Participant is frmaing the experience as a modern version of games she used to play in her childhood.


        \item Hidden secret anticipation
    
        \begin{quote}
            "I expect some surprises and presents in this version for heart-seeking players."
        \end{quote}
        Voice mode presented as a modality which allows for easter eggs to be included for boundary-seeking players.


        \item Voice as novelty
    
        \begin{quote}
            "It's something new and it's cool to see something new... I think it would be even cooler if you could do more things." 
        \end{quote}
        Appreciation for the concept despite it's current limitations or stressors caused by it's form.


        \item Gut-feeling immersion flip
    
        \begin{quote}
            "The first version, I think... I actually don't know how to justify that, but this is just a gut feeling." 
        \end{quote}
        Despite initial answer reverse immersion answer on probing. The participant can't quite explain herself but attributes it to a strong gut-feeling.


        \item Future experience optimism
    
        \begin{quote}
            "Now if I was to play another time, I would probably be more focused on really choosing the plan... a second take would be much easier."
        \end{quote}
        Participant showing that she believes that her answers wouldn't been different had she had more time to explore the game. 


        \item Descripions importance
    
        \begin{quote}
            "I have to pay attention to what's written here... Let's check the boxes. Is there anything inside that can be useful?" 
        \end{quote}
        The player analysed the descriptions on a very deep level. She looked for clues, she read every word. Importance of the description fields shows up here.
         

\end{enumerate}


\subsection{Searching for Themes}

\subsubsection{A - Cooperative nature}
Just like participant \#2 Klara was phrasing her commands in a cooperative manner by using the "Let's.." phrasing. 


\subsubsection{B - Discovery tools}

Klara noted that voice can be used as a discovery tool that would to explore the game space in a way that the text interaction wouldn't make possible. 


\subsubsection{C - Immersion as a gut feeling}

The player described the keyboard version as more immersive and then did a 180 switch based on her gut feeling. Her model of immersion seems to be highly unusual further adding to how differently immersion seems to work across tested participants.


\subsubsection{D - Descripion as atmothsere fuel}

the participant noted multiple times how the descriptions made her feel in the game. She viewed the descriptions as something positive which stands in contrast to the other participants.


\subsection{Reviewing themes}


%     Check if those found themes accurately represent the difference between two input modalities, 
%     and if they fit the main research question.


Below for quick reference are the research questions from the thesis: 

\begin{enumerate}
    \item How would the participants describe the experience of using voice input in the presented
game environment?
    \item How do users feel about using the voice modality compared to traditional input?
\end{enumerate}

The themes are anylysed to see if they fit the main research questions and if they accurately represent the difference between the two input modalities.

\subsubsection{A - Cooperative nature}
% Just like participant \#2 Klara was phrasing her commands in a cooperative manner by using the "Let's.." phrasing. 
This theme directly ansewrs the research questions since the participant is taling about how she's feeling when using the voice commands and the structure of said commands has a direct impact on her emotions and experience.

\subsubsection{B - Discovery tools}

% Klara noted that voice can be used as a discovery tool that would to explore the game space in a way that the text interaction wouldn't make possible. 
Using voice as an exploratory tool is a topic about the player experience and comparing it to the keyboard version. This is directly related to the research questions.


\subsubsection{C - Immersion as a gut feeling}

% The player described the keyboard version as more immersive and then did a 180 switch based on her gut feeling. Her model of immersion seems to be highly unusual futher adding to how differently immersion seems to work across tested participants.

Immersion is a feeling, a feeling that's highly personal and highly impactful when playing any game. By discussing immersion the research problems are directly tackled. 


\subsubsection{D - Descripion as atmothsere fuel}

% the participant noted multiple times how the descriptions made her feel in the game. She viewed the descriptions as something positive which stands in contrast to the other participants.

The descriptions and their impact on the atmosphere is also connected to the research question - the role that they play in both version is different, therefore they impact the player experience during the playhrough.




\subsection{Refining Themes}

\subsubsection{A - Cooperative nature}
% Just like participant \#2 Klara was phrasing her commands in a cooperative manner by using the "Let's.." phrasing. 
This theme directly ansewrs the research questions since the participant is taling about how she's feeling when using the voice commands and the structure of said commands has a direct impact on her emotions and experience.

The participant is framing her interaction as a cooperative endeavour. She's using the "Let's" phrasing showcasing how for her this is an experience in which she's not alone but with the computer together. She's not addressing the game as a servant or a medium that she's asking for permission. Rather, she directly asks the computer if they can do things together.

\subsubsection{B - Discovery tools}

% Klara noted that voice can be used as a discovery tool that would to explore the game space in a way that the text interaction wouldn't make possible. 
Klara envisioned the voice modality as a tool that can be used for slow, gradual exploration of the game world. She stated that had the been equipped to answer her questions it would allow for a very deep sense of exploration. 
She expressed that she would love to see "Easter Eggs" in the game. Small secrets that are out there waiting to be discovered by the players. A game which has a large amount of data about it's objects that can use it to answer to the player could allow for interesting exploration paths indeed. A player could ask about relationships between objects. This is as much of design description as much as experience report. It specifies not only how does the voice mode feel but also how it could create the feeling of being worthy of the extra effort needed to use this input modality compared to traditional input. Implication here being that the voice-driven version without those extra features is just a more exhausting version of the keyboard navigation experience. When told that the game contains some of those hidden mechanics that she mentioned she confirms that it was precisely what she hoped for. 

This is interesting since she's the only player who actively probed the system in this direction. 

\subsubsection{C - Immersion as a gut feeling}

% The player described the keyboard version as more immersive and then did a 180 switch based on her gut feeling. Her model of immersion seems to be highly unusual futher adding to how differently immersion seems to work across tested participants.

Klara gave two conflicting statements on the manner of the game immersion. At first she said that keyboard version is more immersive, only to reverse course and say that the voice-driven version offered deeper immersion for her. What stands out is that fact the she couldn't quite explain this sudden change quoting her gut feeling. This could showcase that at least for her the feeling of immersion has two layers, one of them being the explainable awareness, and the other being on a deeper layer of her brain. This is pure speculation of course, but it seems to explain her personal behaviour.

This complies with the rest of the participants and they way that they all tend to see immersion differently.



\subsubsection{D - Descripion as atmothsere fuel}

% the participant noted multiple times how the descriptions made her feel in the game. She viewed the descriptions as something positive which stands in contrast to the other participants.

The descriptions and their impact on the atmosphere is also connected to the research question - the role that they play in both version is different, therefore they impact the player experience during the playhrough.

The participant relationship with the location descriptions seemed to have been more literary than purely functional. She made multiple remarks on the quality of the writing that went into the descriptions. She made pauses during the gameplay to read each location desciption carefully, often remarking out loud on how she found them "romantic" or "poetic". This type of exploration seemed to have been forced by the voice navigation framework, since she had no other way but to read the descriptions to be able to navigate - especially since she lacked the context clues from the baseline version. Oddly enough she didn't stop her behaviour even when playing the keyboard version suggesting that for her the descriptions played a major importance in the game. At the same time it's worth noting that she's the only participant who volountaryliy read all of the descriptions. All other participants either remarked that the descriptions were too long / unneeded or they found ways around them to be able to progress without having to read them. 

This state of forced immersion in the worm of descriptions is an interesting effect, but in this case it wasn't even needed. Klara simply enjoyed the text as is, even if it was entirely optional for her success. This goes hand in hand with traditional games in which there is often a layer of text in form of collectible items or Easter Egggs that deepens the immersion or expands on the game narrative. Only a subset of players actively engage with this optional content. Klara seems to be just that type of a player.


\subsection{Writing up}


\subsubsection{A - Cooperative nature}
% Just like participant \#2 Klara was phrasing her commands in a cooperative manner by using the "Let's.." phrasing. 
This theme directly ansewrs the research questions since the participant is taling about how she's feeling when using the voice commands and the structure of said commands has a direct impact on her emotions and experience.

\begin{quote}
    "I never even thought about it. It was just natural. 'Let's go here, let's go there, I want to do this, I want to do that' — that was obvious to me. I feel like we are playing together, me and the computer, and not like I'm asking him to take me somewhere. It's like we are together in this."
\end{quote}

The participant is framing her interaction as a cooperative endeavour. She's using the "Let's" phrasing showcasing how for her this is an experience in which she's not alone but with the computer together. She's not addressing the game as a servant or a medium that she's asking for permission. Rather, she directly asks the computer if they can do things together.

She also framed the voice interaction as the "AI" era modality compared to text based input:
\begin{quote}
    "The first version makes me think more of a ChatGPT stuff and AI conversation and the other version makes me think of those old 80s games from Sierra... The first version is more up-to-date because it makes me think it's more connected to the AI and the fact that you can talk to the computer and it understands you."
\end{quote}

\subsubsection{B - Discovery tools}

% Klara noted that voice can be used as a discovery tool that would to explore the game space in a way that the text interaction wouldn't make possible. 
Klara envisioned the voice modality as a tool that can be used for slow, gradual exploration of the game world. She stated that had the been equipped to answer her questions it would allow for a very deep sense of exploration.

\begin{quote}
    "I would like to check the game's limits. That would be my biggest fun. And I suppose that the game creator could be also very inventive and hide some crazy solutions for those people who want to play this way — say, okay, let's just go to the roof, and then you go to the roof of the building."
\end{quote}

She expressed that she would love to see "Easter Eggs" in the game. Small secrets that are out there waiting to be discovered by the players. 
\begin{quote}
    "I expect some surprises and presents in this version for heart-seeking players — players who are patient and want to explore more than just go through the game quickly. It gives me more opportunities and then I can really, okay, I can even ask him to go and pick this, this, this and check a few places around the room even. Not just to go up and down but let's say let's check this box, let's check this box, let's see if there's anything on the floor."
\end{quote}

A game which has a large amount of data about it's objects that can use it to answer to the player could allow for interesting exploration paths indeed. A player could ask about relationships between objects. This is as much of design description as much as experience report. It specifies not only how does the voice mode feel but also how it could create the feeling of being worthy of the extra effort needed to use this input modality compared to traditional input. Implication here being that the voice-driven version without those extra features is just a more exhausting version of the keyboard navigation experience. When told that the game contains some of those hidden mechanics that she mentioned she confirms that it was precisely what she hoped for. 

\begin{quote}
    "In the first version I have no idea about the limits. I don't know if it's going to put me on the moon if I want to or not. In the second version, I know that I only have four options to choose from, and that's mostly that."
\end{quote}

This is interesting since she's the only player who actively probed the system in this direction. 

\subsubsection{C - Immersion as a gut feeling}

% The player described the keyboard version as more immersive and then did a 180 switch based on her gut feeling. Her model of immersion seems to be highly unusual futher adding to how differently immersion seems to work across tested participants.

Klara gave two conflicting statements on the manner of the game immersion. At first she said that keyboard version is more immersive:

\begin{quote}
    "In the second version... maybe it's also because I already got to know the game a little bit... the second part was easier not only because there were already options waiting for me but also because I already knew the game a little bit."
\end{quote}

But after the researched probed her to elaborate she reversed course and say that the voice-driven version offered deeper immersion for her. What stands out is that fact the she couldn't quite explain this sudden change quoting her gut feeling.

\begin{quote}
    "The first version, I think. The first version, because... because... because... I actually don't know how to justify that, but this is just a gut feeling that the first version."
\end{quote}

This could showcase that at least for her the feeling of immersion has two layers, one of them being the explainable awareness, and the other being on a deeper layer of her brain. This is pure speculation of course, but it seems to explain her personal behaviour.



This complies with the rest of the participants and they way that they all tend to see immersion differently.



\subsubsection{D - Descripion as atmothsere fuel}

% the participant noted multiple times how the descriptions made her feel in the game. She viewed the descriptions as something positive which stands in contrast to the other participants.

The descriptions and their impact on the atmosphere is also connected to the research question - the role that they play in both version is different, therefore they impact the player experience during the playhrough.

The participant relationship with the location descriptions seemed to have been more literary than purely functional. She made multiple remarks on the quality of the writing that went into the descriptions. She made pauses during the gameplay to read each location desciption carefully, often remarking out loud on how she found them "romantic" or "poetic". 

\begin{quote}
    "I think I very much like the descriptions. They are absolutely amazing and very... very romantic I'd say even, perfect description. I would like to see, if I was to play longer, a bit more of the picture because it made me somehow suffocate that I could not see the whole picture."
\end{quote}

This type of exploration seemed to have been forced by the voice navigation framework, since she had no other way but to read the descriptions to be able to navigate - especially since she lacked the context clues from the baseline version. Oddly enough she didn't stop her behaviour even when playing the keyboard version suggesting that for her the descriptions played a major importance in the game. At the same time it's worth noting that she's the only participant who volountaryliy read all of the descriptions. All other participants either remarked that the descriptions were too long / unneeded or they found ways around them to be able to progress without having to read them.  

But she looked for system boundaries by scrutinising the descriptions:
\begin{quote}
    "The description says that someone has been running through this room in a hurry. Look around if there is any door in the room. Can we move forward? Because I would like to use a crowbar for something... I'm trying to catch the logic of it."
\end{quote}

This state of forced immersion in the worm of descriptions is an interesting effect, but in this case it wasn't even needed. Klara simply enjoyed the text as is, even if it was entirely optional for her success. This goes hand in hand with traditional games in which there is often a layer of text in form of collectible items or Easter Egggs that deepens the immersion or expands on the game narrative. Only a subset of players actively engage with this optional content. Klara seems to be just that type of a player.



\section{Language fluency consideration}

% Include the level of English language that the participant used and their percieved degree of fluency.
% Are they native speakers? What about their accent? Make sure to include that too.

The player is on a self proclaimed B2/C1 level in English. It's her second language. She has a subtle Polish accent

\section{Transcript}

Gabriel: okay this is just a start uh all right so I'm sending the recording very good um so the the way this is gonna go down uh first let me just ask you a few simple questions um so first question as I remember is uh what is your gaming expertise so are you a pro gamer are you a casual gamer I know you play the sims

Wicia: I'm a casual gamer, I don't play that much often, and I also play like two games at once, so I wouldn't say I'm a pro gamer.

Gabriel: I mean you play two at once, most people only play one game at a time, so I don't know, I don't know, you tell me.

Gabriel: Okay, and what is your familiarity with voice assistants, like Siri, Google, Alexa?

Wicia: I know of them, but I don't use them at all.

Gabriel: Okay, that's great.

Wicia: I think I'm not comfortable with talking to my phone.

Gabriel: Okay, we'll get back to it then.

Gabriel: And just for the record, if you want to share, from the accessibility point of view, do you have any accessibility constraints like vision impairments, mobility impairments, stuff like that?

Wicia: No, not any I'm aware of.

Gabriel: Okay, that's very good.

Gabriel: Don't worry, this vlog will not test you if you have them.

Gabriel: okay so let me just give you like a brief intro to this this is what you're gonna be playing is a simple video game that I created and the body of my work is about testing the different control modality which is voice based so this game right here I will give you 10 minutes to play it and

Gabriel: if there's any problems or whatever just ask me as we go along, right?

Gabriel: But the idea is that you control the game with your voice and you don't have to express it with direct commands like do this or do that or I want to do this.

Gabriel: You can speak using a natural voice and the game should be able to comprehend and understand you.

Gabriel: of course sometimes things break and then we'll like navigate, circumnavigate it but the game itself is a text-based game so that's the UI, you have the location, oh and

Gabriel: Let me give you a brief story intro, 32nd version.

Gabriel: So it's a post-apocalypse Warsaw and you have a brother and you both, he went for a food run to a shopping mall.

Gabriel: That was two days ago, he hadn't returned.

Gabriel: So you feel like, okay, he's there somewhere, probably like something happened to him and you try and like either confirm he's dead or help him, yeah?

Gabriel: So the navigation is right now purely voice-based.

Gabriel: The way you talk, you are,

Gabriel: press this button or yeah you hold it you say what you want to say and then you release the button and then you wait a little bit because it's gonna be like here like tell you it's processing and then you're gonna give you a result okay all right so again if you have any questions let me know

Gabriel: And if things will be broken, just for the record of this transcription, I will have to restart things.

Gabriel: Okay, so let me start.

Gabriel: And I'll be taking small notes, but I will not be judging how you play.

Gabriel: Don't worry.

Gabriel: And you don't have to finish the game in 10 minutes.

Gabriel: It's not about that.

Gabriel: It's also not about testing the game itself.

Gabriel: It's more about testing the interaction layer, let's say.

Gabriel: Okay?

Wicia: Okay.

Gabriel: Fantastic.

Gabriel: So, please, you can start.

Gabriel: i'll just start a timer in the background yeah take the computer i mean okay

Gabriel: Do you have a pen?

Gabriel: Excuse me to ask now, but I don't think I have that.

Wicia: I have a few of them actually.

Gabriel: Very good.

Gabriel: Thank you.

Gabriel: Alright.

Wicia: Okay, so I can just tell you anything I want to do right now?

Gabriel: Yeah, yeah, yeah, yeah.

Gabriel: hold this button right here and then you say something like oh oh shit i clicked it so you have to wait a few seconds uh but you hold it and then you say like i wanna maybe travel there or i want to do this pick this up yeah i want to enter the clothing store

Wicia: okay so it's like a chocolate something like that uh maybe a little bit we'll talk about it but don't worry okay

Wicia: I want to open the door with crowbar I don't think

Gabriel: This is the one part I have not explained.

Gabriel: Here you have your inventory.

Gabriel: Here are the items that you have currently empty.

Gabriel: Yeah, so yeah, you don't have a crowbar in your inventory.

Wicia: Okay.

Wicia: Can I go through the ground entrance?

Gabriel: Yeah.

Wicia: I want to enter via ground entrance.

Gabriel: Typically, you're gonna be having the

Gabriel: wall of text of like what you can do yeah and at the end are the possible locations you can travel to okay so that just to like make it a little bit clearer can i pick items like around me can i like go to the clothing store and pick the shirt and something like that technically yeah yeah yeah okay also take note that you can go to another place as well which is the this place right here yeah i know yeah uh

Wicia: but yeah you okay yeah please i want to enter the small storage shed

Wicia: Crowbar.

Wicia: You got it, you got it.

Wicia: I want to pick up the crowbar I want to leave this place

Wicia: I want to enter the clothing store.

Wicia: And now I want to open the door.

Wicia: Do I have to say that I want to open the door, the thing, or can I just say that?

Gabriel: We'll see, I'm not sure.

Gabriel: As far as I'm, it really depends.

Wicia: Oh, okay.

Wicia: I want to enter the unlocked room.

Oh, okay.

Okay, I want to enter a clothing store.

Wicia: i want to enter the corridor hey

Wicia: I want to enter the fishing store.

Gabriel: The fishing store has no graphic, but it's just very dark here.

Wicia: Yeah.

Gabriel: okay so i cannot use anything if it's not mentioned in the text you mean yeah you you can try right i'm not saying you can't but the text should be a guidance here right yeah i would say it's rather improbable that you can use i think there are a few cases but only a few cases

Wicia: Okay, I want to leave the fishing store.

Wicia: I want to enter the food court.

Wicia: can I enter the bubble tea store?

Wicia: go for it I want to enter a bubble tea store yay

Just for the record, we'll have to leave at what time? 17.20?

Wicia: Yeah, I think something like that.

Gabriel: Okay, okay, okay.

Gabriel: Well, then we're gonna do this part two minutes shorter than usual.

Gabriel: Don't worry about that.

Gabriel: So, right now I will do one thing, which is to do this.

Gabriel: And now I want you to play for another eight minutes, maybe six minutes, using the keyboard.

Gabriel: So now, no talking.

Gabriel: you use the keyboard here to navigate arrows and enter sometimes when you like use items you select them from here okay so i see that i don't have any other options than leave this place yeah but i see that there's an action crafting yeah with crafting you can combine two items to get another item

Wicia: ♪ Shadow me through all the hurt and pain ♪ ♪ Let the heart do the talking ♪ ♪ When you have nothing left to say ♪

Wicia: ♪ Church bells, do you still remember?

♪ ♪ Oh, let me be your housekeeper, yeah ♪ ♪ Come to me, come to me ♪ ♪ Oh, let me be your housekeeper, yeah ♪ ♪ Come to me, come to me ♪

Gabriel: here there's a small arrow so you use the arrows here then return enter to try and use it yeah

Wicia: Oh, OK. You can't, yeah, not here.

Wicia: OK.

Yeah.

Yeah.

Oh.

oh okay

th th

Wicia: I'm going to try and get out of here as quickly as possible.

and an

Wicia: Can I click 4 once again if it didn't... Yeah, that is a bug actually.

Gabriel: Technically you could pick up 10 of them if you want to, or 20.

Wicia: Okay, I see.

Gabriel: It's a very big gun store, okay?

Wicia: That's what I thought at first, but I see that it should be 1.

Gabriel: It should be 1, technically yes.

Gabriel: The last person I tested with picked up 10 of them, without saying anything.

Gabriel: They were very happy about it.

th th

Gabriel: That's a bug.

Gabriel: Let's assume that this works, okay?

Gabriel: Okay.

Gabriel: Yeah, he's with you.

Gabriel: Just the one thing, because of this bug, you cannot leave the building the same way you came.

Gabriel: So just the way you interact with the game, you cannot just leave through the main door.

Gabriel: okay yeah you have to figure out some other way of quitting the building let's say okay

Wicia: get broken

Wicia: can I have any hints from you?

Gabriel: I mean technically to be fair we are running out of time so let's just assume that you can leave the building so like you know how to do this so technically like you beat the game let's say okay not that it matters a lot in this context but

Gabriel: Congratulations, I think you did it.

Wicia: Thanks.

Gabriel: Yeah, yeah, yeah, so a pro gamer after all.

Gabriel: Okay, so I know that we have to go, so I propose that I will have to ask you a few questions, so I can ask them as we walk.

Gabriel: If that's okay, I'll just have to send them to my phone.

Gabriel: Yeah, because we have to get going, right?

Wicia: Do you think we can leave 6 minutes?

Gabriel: Ok, then maybe we can.

Gabriel: Let me just ask you a few simple questions.

Gabriel: Now that you tried and replayed both of these versions, how are you feeling and was there anything that immediately surprised you?

Wicia: okay so um i think it was easier to play with like clicking buttons than speaking but it's because i'm not comfortable speaking i'm not used to it and i think that

Wicia: It's a pretty creative way to play games.

Wicia: I haven't played any games like this before, so it was something new for me.

Wicia: And I think it gives you the illusion of choice, because you don't know what you can do and what you cannot do.

Wicia: Yeah, I think it was great.

Okay.

Gabriel: And if you had to describe the difference between the two games to a person that hasn't played the game at all, how would you describe the difference?

Wicia: Okay, so I would say that the second version with clicking is more simple and can be much faster because you don't have to think what you can do and what you cannot.

Wicia: You have simple choices that you can see.

Wicia: And with talking you have to use your imagination with what you can do and what you want to do.

Wicia: So I would say its second version is more limited to something.

Wicia: Yeah, I think I would explain it that way.

Gabriel: From the moment that you used your voice to navigate, what were you thinking as you decided what to say to the game?

Wicia: I didn't know how to pick the best words because I thought to myself, will it understand me?

Wicia: So I was trying to be as clear as possible and as official as I can be.

Gabriel: Yeah, I understand you.

Gabriel: Yeah.

Gabriel: And how was it like to use the voice input in the game?

Gabriel: How was it like to use your voice to navigate?

Wicia: It was weird.

Wicia: As I said, it's because I don't use my voice to anything else than communicate with others, with other people.

Wicia: I do not use anything that requires voice guidance.

Wicia: Yo, yo, yo.

Wicia: So I think It was weird for me, but it will be easier for people that use voice comments to communicate via Alexa or Google Assistant You don't talk to chat GPT for instance, do you?

Gabriel: No That's okay, no worries

Gabriel: In which version did you feel more inside the game, more inside the story world?

Wicia: I think in the second one, because I could look at the screen and don't have to think too much of what to pick.

Wicia: and how to create proper sentence and I think it's easier for me to pick options that are already available and not to give comments on my own because I

Wicia: I think if I had an idea of what I would be playing, then I would be more comfortable with talking, but because it was like told me on the spot, I couldn't prepare myself and I think the second option was more comfortable for me.

Gabriel: And how did the effort of thinking what to say compare to the effort of just pressing, knowing which key to press?

Wicia: I think that I read much faster than I speak.

Wicia: I think I was doing much much better in the second option because I was just sometimes I didn't finish reading and I could pick an option for finishing reading so it was easier for me and I think I saved some time before because of that and I did better progress in the second time

Gabriel: I understand.

Gabriel: And did the freedom, you know, to express anything, to say anything to the game, did it make you feel more in control or was it more overwhelming than the second version?

Wicia: I think it was overwhelming, but it's because English is not my first language.

Wicia: I think it will be easier when I would

Wicia: be able to speak Polish and because I'm stressed because we're in public.

Wicia: But I think if I was in like more comfortable place for me then it would be easier for me and it wouldn't be as stressful as it was.

Gabriel: I understand.

Gabriel: And how would you compare the overall experience of these two versions?

Gabriel: Which one would you choose and why?

Wicia: I think

Wicia: On the first thought, I would choose the second one.

Wicia: But when I think about it more, I think the first option would be more challenging for me.

Wicia: And if I could test it, I don't know, in my home, in some more quiet environment, that it would be a good way to practice speaking to the machines, I guess.

Gabriel: Yeah, yeah.

Gabriel: I understand.

Wicia: I think it's a good way to be more comfortable with speaking.

Wicia: If you have any problems with communicating with others and you are worried of talking to other people, I think it's a good exercise for that.

Wicia: Because you need to think and properly... How to say that?

Wicia: Exonerate?

Wicia: I think I would say that.

Wicia: Your thoughts.

Gabriel: In the beginning you said that it is kind of like a chatbot.

Gabriel: What do you mean by that?

Wicia: When I was in the high school, there was this chatbot game when you picked some options and then the game, I mean the chatbot, the chatbot will show you the plot.

Wicia: You can just choose something and it will

Wicia: walk you through their plot.

Wicia: I don't know how to explain.

Wicia: I can't even remember what the game is called.

Wicia: But it was something like that.

Wicia: So I've played something similar and it was Childhood Mechanics.

Wicia: I think it was on Nintendo.

Wicia: Okay, I see, I see.

Gabriel: Okay, interesting.

Gabriel: Is there anything else that you would like to add or ask me?

Wicia: I think it's something new and it's cool to see something new.

Wicia: It's a pretty chill game to play and it's cool that you can interact with every room that you enter and you can take a look around and see what's around you.

Wicia: Yeah, I think it would be even cooler if you could do more things and if you could try to do more stuff.

Wicia: But overall, I think it's something new and something that can be revolutionary.

Gabriel: Okay, thank you very much for your input.

Gabriel: This is tremendously helpful.

Gabriel: And this is the end of this recording.






